\section{Gaussian comparison of ordinary fluctuations}
\label{gs:section}

This section proves a reference-law estimate for the \emph{global} ordinary
quadratic source when all coordinates on a thin separator are observed. The
proof retains the cross-block contractions, arbitrary signs of the original
marks, and the mean across commuting flat holonomies. It also shows why the
second-moment gain cannot by itself be promoted to a uniform higher-source
bound after random-offset averaging. None of these statements replaces the
actual compact marginal by a Gaussian law.

\subsection{The original Maxwell quadratic and its positive root matrices}

Work at unit mesh on the periodic four-dimensional torus of side \(n\ge3\).
Put \(V=n^4\) and \(d_G=\dim SU(N)\). On momentum
\(k\in(2\pi/n)\{0,\ldots,n-1\}^4\), write
\[
 p_\mu=e^{ik_\mu}-1,\qquad u_\mu=|p_\mu|^2,\qquad
 s=\sum_\mu u_\mu,\qquad K=sI-pp^*.
\]
Use the translation-equivariant Coulomb section. After quotienting gradients,
omit only the remaining constant harmonic zero modes of this specified
reference. All traces below are real Hilbert-space traces; conjugate Fourier
and color sectors are paired with their actual real multiplicities.
Let \(K^{+1/2}=(K^+)^{1/2}\), extended by zero on the kernel. Retain the original
quadratic normalization \(\upsilon>0\):
\begin{equation}
 S_2(A)=\frac{\upsilon}{2}\inner A{KA},\qquad
 A=(\beta\upsilon)^{-1/2}K^{+1/2}\xi,
 \label{gs:reference}
\end{equation}
where \(\xi\) is standard real Gaussian on the positive transverse space.

The directional row with root direction \(\nu\) and plane partner \(\mu\ne\nu\)
is
\[
 L_{\nu\mu}A=\overline p_\mu(p_\mu A_\nu-p_\nu A_\mu),
 \qquad L_{\nu\mu}p=0.
\]
Its original linear force is \(\upsilon L_{\nu\mu}A\). Equivalent fixed
edge-placement conventions introduce unit phases and preserve the following
norms and traces. For each literal oriented-link root \(c=(x,\nu)\), let
\(T_c\) sum its three row squares and its color components, with the unitary
Fourier normalization \(V^{-1/2}\). Set
\begin{equation}
 B_c=K^{+1/2}T_cK^{+1/2},\quad B_f=\sum_cf_cB_c,\quad
 B_{\rm all}=\sum_cB_c,\qquad
 Q_f=\upsilon\xi^*B_f\xi-\frac13\ell^{(0)}[f].
 \label{gs:source}
\end{equation}
The deterministic Casimir scalar is retained. This is the leading complete
ordinary quadratic in the named Maxwell reference. Lower field-degree-two
corrections are addressed separately below, rather than silently identified
with \eqref{gs:source}.

\begin{lemma}[Complete root bounds]
\label{gs:rootbounds}
For the literal original oriented roots,
\begin{equation}
 B_c\ge0,\qquad \frac13K\le B_{\rm all}\le8I,
 \qquad \Tr B_c\le6d_G.
 \label{gs:bounds}
\end{equation}
The operator inequalities are on the positive transverse space.
\end{lemma}
\begin{proof}
Put \(C_{\mu\nu}=p_\mu A_\nu-p_\nu A_\mu\). Directly,
\[
 \inner A{KA}=\sum_{\mu<\nu}|C_{\mu\nu}|^2,\qquad
 \sum_{\nu,\mu\ne\nu}|L_{\nu\mu}A|^2
 =\sum_{\mu<\nu}(u_\mu+u_\nu)|C_{\mu\nu}|^2.
\]
Since \(u_\mu\le4\), \(\sum_cT_c\le8K\). Conversely
\((KA)_\nu=\sum_{\mu\ne\nu}L_{\nu\mu}A\), so Cauchy over its three partners
gives \(K^2\le3\sum_cT_c\). Conjugate both inequalities by \(K^{+1/2}\).
For the root trace, the rows annihilate gradients and \(K=sI\) transversely;
therefore
\begin{equation}
 \Tr B_{(x,\nu)}
 =\frac{d_G}{V}\sum_{k:s(k)>0}\sum_{\mu\ne\nu}
       \frac{u_\mu(u_\mu+u_\nu)}s
 \le \frac{d_G}{V}\sum_k\sum_{\mu\ne\nu}u_\mu
 =6d_G.
 \label{gs:roottrace}
\end{equation}
The last equality uses the exact one-coordinate grid average \(\E_k u_\mu=2\).
The omitted zero mode has zero source amplitude. No inverse lower spectral
bound is used.
\end{proof}

Every fixed pointing/completion of the prescribed source keeps its complete
root list. Its bounded original incidence changes constants by that declared
factor; it does not authorize deleting neighboring or retained roots.

\subsection{All separator coordinates and the full trace kernel}

Let \(T\) be independent of the Gaussian field and let
\(Y_T=L_T\xi\) record \emph{all} prescribed separator link coordinates in the
same section. In particular \(L_T\) includes the covariance square root in
\eqref{gs:reference}; this is not observation of independent local kinetic
pieces. With the Moore--Penrose inverse define
\[
 P_T=L_T^*(L_TL_T^*)^+L_T.
\]
It is the orthogonal projection onto \(\operatorname{ran}L_T^*\), regardless of
linear redundancies among the observations. Exact Gaussian conditioning gives
\begin{align}
 \E(Q_f\mid T,Y_T)
 &=\upsilon\{\xi^*P_TB_fP_T\xi+\Tr[(I-P_T)B_f]\}
       -\ell^{(0)}[f]/3,\label{gs:conditional}\\
 G_{T,f}&:=\E(Q_f\mid T,Y_T)-\E Q_f
   =\upsilon\{\xi^*C_{T,f}\xi-\Tr C_{T,f}\},\quad
 C_{T,f}=P_TB_fP_T.\label{gs:centered}
\end{align}
For each \(T\), \(\E_\xi G_{T,f}=0\). Thus the expectation of the original
source does not fluctuate with the independent offset. Its exact conditional
variance is
\begin{equation}
 \Var_\xi G_{T,f}=2\upsilon^2\Tr(C_{T,f}^2).
 \label{gs:exactvariance}
\end{equation}

\begin{theorem}[Arbitrary signed marks under translation averaging]
\label{gs:separator}
Suppose the observation is a uniform translation of one fixed separator pattern
in the translation-equivariant reference, and \(\operatorname{rank}P_T\le r\).
For every real deterministic original mark \(f\),
\begin{equation}
 \Var_{T,\xi}\E(Q_f\mid T,Y_T)
 \le128\upsilon^2\frac rV M_f,
 \qquad M_f=\sum_cf_c^2.
 \label{gs:rankgain}
\end{equation}
In particular a separator containing at most \(CV/R\) links obeys
\begin{equation}
 \Var_{T,\xi}\E(Q_f\mid T,Y_T)
 \le C\upsilon^2d_G R^{-1}M_f.
 \label{gs:separatorgain}
\end{equation}
\end{theorem}
\begin{proof}
The full root covariance kernel is
\[
 \mathcal K_T(c,d)
 =2\upsilon^2\Tr(P_TB_cP_TB_dP_T)
 =2\upsilon^2\norm{B_c^{1/2}P_TB_d^{1/2}}_{\HS}^2\ge0.
\]
This retains every contraction between roots in different cubes. Its row sum
uses the \emph{global} positive source bound:
\begin{equation}
 \sum_d\mathcal K_T(c,d)
 \le16\upsilon^2\Tr(P_TB_c).
 \label{gs:rowsum}
\end{equation}
For each orientation separately, translation covariance and uniform averaging
give
\[
 \E_T\Tr(P_TB_{(x,\nu)})
 =\frac1V\sum_y\Tr(P_0B_{(y,\nu)})
 \le\frac1V\Tr(P_0B_{\rm all})\le\frac{8r}{V}.
\]
No rotational symmetry among the orientation types is needed. Positivity,
symmetry and \(2|f_cf_d|\le f_c^2+f_d^2\) now give
\[
 \sum_{c,d}f_cf_d\E_T\mathcal K_T(c,d)
 \le\sum_cf_c^2\sum_d\E_T\mathcal K_T(c,d)
 \le128\upsilon^2(r/V)M_f.
\]
Each observed link supplies at most \(d_G\) real coordinates, so the rank is
at most \(Cd_GV/R\). Equation~\eqref{gs:centered} ensures that the total
variance contains no further between-offset mean term.
\end{proof}

One may average over the full torus orbit of any fixed packing. For a period
\(R\) pattern with \(R\mid n\), averaging its period cell suffices. If a
different offset convention leaves a nonperiodic remainder, the translation
identity must be checked for that geometry. Before averaging, the proof gives
the useful but weaker spatially weighted inequality
\begin{equation}
 \Var_\xi G_{T,f}
 \le16\upsilon^2\sum_cf_c^2\Tr(P_TB_c).
 \label{gs:fixedoffset}
\end{equation}
Rank alone does not give the \(R^{-1}\) gain for arbitrary marks at each fixed
offset.

\subsection{Sharpness for a full separator}

\begin{proposition}[The \(R^{-1}\) power cannot be improved in general]
\label{gs:sharpness}
Suppose \(R\mid n\) and the observed separator contains all tangential links
on the \(n/R\) hyperplanes \(x_1=T_1\pmod R\). For the global source
\(f_c=1\), every offset satisfies
\begin{equation}
 \Var_\xi\E(Q_1\mid T,Y_T)
 \ge c\upsilon^2d_G V/R.
 \label{gs:sharp}
\end{equation}
Since \(M_1=4V\), this matches \eqref{gs:separatorgain} in its power of \(R\).
\end{proposition}
\begin{proof}
Consider real link covectors \(h\) supported on the indicated hyperplanes,
with \(h_1=0\). At each tangential momentum \(k'=(k_2,k_3,k_4)\), require
\[
 \sum_{j=2}^4\overline p_jh_j=0,
 \qquad s'(k')=\sum_{j=2}^4|p_j|^2\ge1.
\]
There are at least \(cn^3\) such momenta and two tangential transverse
polarizations at each, for each hyperplane and color. Pairing conjugate modes
gives a real space \(\mathcal H\) of dimension at least \(cd_GV/R\).
Every member is globally divergence-free and has no harmonic zero component.
It is supported on observed coordinates. Consequently
\[
 \mathcal E=K^{+1/2}\mathcal H\subseteq\operatorname{ran}P_T,
 \qquad \dim\mathcal E=\dim\mathcal H.
\]
The inverse square root is injective there and preserves the tangential
frequency support. On \(\mathcal E\), \(K\ge I\), hence
\(B_{\rm all}\ge K/3\ge I/3\) as a quadratic-form bound. If \(\Pi\) projects
onto \(\mathcal E\), Hilbert--Schmidt contraction under compression gives
\[
 \Tr[(P_TB_{\rm all}P_T)^2]
 \ge\Tr[(\Pi B_{\rm all}\Pi)^2]\ge\dim\mathcal E/9.
\]
Apply \eqref{gs:exactvariance}. Additional observations enlarge the observed
subspace and preserve this lower bound.
\end{proof}

The full separator retains high tangential frequencies. A stronger estimate
for a separate low-mode observation, such as an \(R^{-8}\) estimate in the
parallel infrared reference, cannot be substituted here. The proposition is
about the complete global quadratic; it does not tile a physical-scale
nonlinear induction.

\subsection{Signed quadratic jets with a positive matrix envelope}

The positivity of each \(B_c\) can be replaced by a quantitative positive
envelope. This allows an explicitly controlled lower-derivative quadratic
partner to remain in the reference observable.

\begin{proposition}[Envelope form of the separator estimate]
\label{gs:envelope}
Let \(A_c\) be real symmetric matrices, covariant under translations with a
fixed finite set of root types per cell. Suppose positive matrices \(C_c\)
satisfy
\[
 -C_c\le A_c\le C_c,\qquad \sum_cC_c\le\Lambda I.
\]
Let \(V_{\rm cell}\) denote the number of translation cells. For
\(H_f=\xi^*(\sum_cf_cA_c)\xi-\Tr(\sum_cf_cA_c)\), under the same observation,
\begin{equation}
 \Var\E(H_f\mid T,Y_T)
 \le2\Lambda^2\frac r{V_{\rm cell}}\sum_cf_c^2.
 \label{gs:envelopebound}
\end{equation}
\end{proposition}
\begin{proof}
The Loewner bounds imply that \(A_c\) vanishes on \(\ker C_c\), so
\(A_c=C_c^{1/2}J_cC_c^{1/2}\) with a selfadjoint contraction \(J_c\), even
when \(C_c\) is singular. For \(U=C_c^{1/2}PC_d^{1/2}\), Schatten Cauchy gives
\[
 |\Tr(PA_cPA_dP)|=|\Tr(J_cUJ_dU^*)|
 \le\norm U_{\HS}^2=\Tr(PC_cPC_dP).
\]
Thus the absolute full covariance kernel is dominated by the positive envelope
kernel. Its row sum is at most \(2\Lambda\Tr(PC_c)\). Averaging over translations
gives \(\E_T\Tr(P_TC_c)\le\Lambda r/V_{\rm cell}\), separately for each type.
The signed Schur argument in Theorem~\ref{gs:separator} proves the claim.
\end{proof}

For example, suppose a \emph{specified fixed-cell} compact flat source jet,
evaluated in its named Gaussian marginal, has whitened root matrices
\[
 A_c=\tfrac12\bigl(\mathsf R_c^*\mathsf R_c
        +\beta^{-1}\mathsf D_c^*\mathsf J_{\beta,c}\mathsf D_c\bigr),
\]
with the already established fixed-cell bounds
\[
 \sum_c\mathsf R_c^*\mathsf R_c\le CI,\quad
 \sum_c\mathsf D_c^*\mathsf D_c\le CI,\quad
 \norm{\mathsf J_{\beta,c}}_{\rm op}\le C.
\]
Taking
\(C_c=\tfrac12(\mathsf R_c^*\mathsf R_c+
C\beta^{-1}\mathsf D_c^*\mathsf D_c)\) gives
\(\Lambda\le C(1+\beta^{-1})\). If this same reference observation has rank
\(Cd_GV/R\), its \emph{complete} quadratic jet therefore satisfies
\[
 \Var(\hbox{retained quadratic jet})
 \le Cd_GR^{-1}(1+\beta^{-1})^2M_f.
\]
This is an implication under the displayed fixed-cell envelope, consistent with
the corresponding jet estimate of \texttt{block\_and\_transfer/ordinary\_law/compact\_source\_jets.tex}. It does not supply that
envelope uniformly for a growing eliminated cube. In that geometry the inverse
block, actual source Hessians and their boundary responses require separate
bounds. It also makes no comparison between a Gaussian marginal and the
original compact induced marginal.

\subsection{Uniform flat-holonomy mixtures and their common mean}

Fix a commuting flat background in a constant representative. In an adjoint
color sector let \(\theta=(\theta_1,\ldots,\theta_4)\) be the total torus twist.
Its covariant symbols are
\[
 p_\mu(k,\theta)=\exp\!\left(\frac{i(2\pi k_\mu+\theta_\mu)}n\right)-1,
 \quad u_\mu=|p_\mu|^2,\quad s=\sum_\mu u_\mu,
 \quad K_\theta=sI-pp^*.
\]
Conjugate color sectors have opposite twists and form a real Gaussian pair;
Cartan sectors are untwisted. Use the translation-equivariant covariant Coulomb
section and omit only its true covariant harmonic zero modes. Positive modes
arbitrarily close to zero remain in the reference.

The proofs of \eqref{gs:bounds} and Theorem~\ref{gs:separator} are uniform
in \(\theta\): shifted-grid averaging still gives \(\E_k u_\mu=2\), and the
rank bound only counts the real observed coordinates. In particular,
\begin{equation}
 \E_T\Var_\xi\!\left(\E(Q_f\mid\theta,T,Y_T)\mid\theta,T\right)
 \le C\upsilon^2d_GR^{-1}M_f.
 \label{gs:twistedconditional}
\end{equation}
The remaining issue within this mixture is its \emph{between-holonomy mean};
it must not be discarded by centering separately at each holonomy.

\begin{lemma}[Uniform root-mean variation]
\label{gs:holonomymean}
For one color and root orientation, define on the momentum torus
\[
 b_\nu(k)=\sum_{\mu\ne\nu}
       \frac{u_\mu(k)[u_\mu(k)+u_\nu(k)]}{s(k)},\qquad b_\nu(0)=0.
\]
Its shifted-grid average satisfies
\begin{equation}
 \left|\frac1{n^4}\sum_kb_\nu((2\pi k+\theta)/n)
       -\frac1{n^4}\sum_kb_\nu(2\pi k/n)\right|
 \le Cn^{-6},
 \label{gs:rootmeanvariation}
\end{equation}
uniformly in every twist, including degenerating twists.
\end{lemma}
\begin{proof}
Near the unique singular point modulo \(2\pi\), \(s(k)\asymp|k|^2\) and the
numerator vanishes to order four. Thus, away from that point,
\[
 |\partial^J b_\nu(k)|\le C_J|k|^{2-|J|}.
\]
For a smooth dyadic annulus of radius \(\rho\), the localized symbol has
\(L^1\) norm \(O(\rho^6)\) and its \(m\)-th derivative has \(L^1\) norm
\(O(\rho^{6-m})\). Integrating by parts in a coordinate with
\(|z_j|\ge|z|/2\), for any fixed integer \(m>6\), bounds that annulus's
Fourier coefficient by
\[
 C\rho^6(1+\rho|z|)^{-m}.
\]
Sum separately over \(\rho\le|z|^{-1}\) and \(\rho>|z|^{-1}\); both
geometric sums are \(O(|z|^{-6})\). The smooth portion away from the singular
point decays faster. Therefore
\begin{equation}
 |\widehat b_\nu(z)|\le C(1+|z|)^{-6}.
 \label{gs:fourierdecay}
\end{equation}
These coefficients are absolutely summable in dimension four. The Fourier
series represents the continuous symbol, including its assigned value zero.
The shifted-grid Poisson identity is consequently absolutely justified:
\[
 \frac1{n^4}\sum_kb_\nu((2\pi k+\theta)/n)
 =\sum_{m\in\mathbb Z^4}\widehat b_\nu(nm)e^{im\cdot\theta}.
\]
Subtract the formula at zero twist and sum \(m\ne0\) using
\eqref{gs:fourierdecay}. This proves \eqref{gs:rootmeanvariation}.
\end{proof}

Write \(m_f(\theta)=\E_\xi(Q_f\mid\theta)\). Every plaquette of the flat
background is identity, so the original Casimir scalar is independent of
\(\theta\). Summing the real color sectors and using \(4V\) literal roots gives
\begin{align}
 |m_f(\theta)-m_f(0)|
 &\le C|\upsilon|d_Gn^{-6}\norm f_1
 \le C|\upsilon|d_Gn^{-4}\sqrt{M_f},\label{gs:meancomparison}\\
 \Var_\pi m_f(\theta)&\le C\upsilon^2d_G^2n^{-8}M_f,
 \label{gs:betweenholonomy}
\end{align}
for \emph{any} source-independent probability law \(\pi\) on commuting flat
holonomies. No independence among color-sector means has been assumed.

\begin{proposition}[The complete flat-holonomy mixture]
\label{gs:holonomy}
For any source-independent holonomy law \(\pi\), with \(T\) independent of
\(\theta\) and of the conditional Gaussian field, the complete observed mean
satisfies
\begin{equation}
 \Var\E(Q_f\mid\theta,T,Y_T)
 \le C\upsilon^2\bigl(d_GR^{-1}+d_G^2n^{-8}\bigr)M_f.
 \label{gs:fullholonomy}
\end{equation}
\end{proposition}
\begin{proof}
The centered mean splits exactly as
\[
 \E(Q_f\mid\theta,T,Y_T)-\E Q_f
 =G_{\theta,T,f}+[m_f(\theta)-\E_\pi m_f],
 \qquad \E(G_{\theta,T,f}\mid\theta,T)=0.
\]
The two terms are orthogonal, while the second square remains. Combine
\eqref{gs:twistedconditional} and \eqref{gs:betweenholonomy}.
\end{proof}
A coarser positive observation of this same Gaussian mixture contracts the
complete bound. The estimate does not assert analyticity of the changing
harmonic quotient through a degenerate twist, nor does it remove harmonic
variables from an actual compact history. Higher holonomy/source derivatives
and a comparison with the nonlinear induced law remain independent problems.

\subsection{Higher offset cumulants: an exact limitation}
\label{gs:higher}

Return to \eqref{gs:centered} and set \(b_f=\norm f_\infty\). Conditional on
\(T\), the complete Gaussian quadratic cumulants are
\begin{equation}
 \kappa_k(G_{T,f}\mid T)
 =2^{k-1}(k-1)!\upsilon^k\Tr(C_{T,f}^k),\qquad k\ge2.
 \label{gs:conditionalcumulants}
\end{equation}
The root bounds and the signed Schur estimate give
\begin{equation}
 \norm{C_{T,f}}_{\rm op}\le8b_f,\qquad
 v_T:=\Tr(C_{T,f}^2)\le48d_GM_f,
 \qquad \E_Tv_T\le Cd_GM_f/R.
 \label{gs:vmax}
\end{equation}
Thus the \emph{averaged conditional} cumulants have bounds
\(C_k|\upsilon|^kb_f^{k-2}d_GM_f/R\). They are not the unconditional
cumulants after adjoining the offset. Since the conditional first mean is zero,
the fourth cumulant of the complete joint variable is exactly
\begin{equation}
 \kappa_4(G)
 =48\upsilon^4\E_T\Tr(C_{T,f}^4)
      +3\Var_T(2\upsilon^2v_T).
 \label{gs:fourthmixture}
\end{equation}
The second term can cost \(M_f^2/R\). Offset independence of the first mean
does not eliminate fluctuations of higher conditional source coefficients.

\begin{proposition}[No blanket higher-source upgrade]
\label{gs:counterexample}
The positivity, bounded root trace, uniform translation, and rank assumptions
of Theorem~\ref{gs:separator} do not imply a fourth-cumulant bound
\(C M_f/R\), or a Bernstein bound with variance parameter \(CM_f/R\) and a
fixed radius \(|t|b_f\le\delta\), for the complete offset mixture.
\end{proposition}
\begin{proof}
Take \(RM\) independent real standard Gaussian coordinates on a cyclic set.
Let \(S_t\) be its residue classes modulo \(R\), let \(T\) be uniform on
\(\{0,\ldots,R-1\}\), and observe all coordinates in \(S_T\). Set
\(B_c=e_ce_c^*\), so \(\sum_cB_c=I\), \(\Tr B_c=1\), and every observation
has rank \(M\). This family is translation-covariant. Mark only \(S_0\):
\(f_c=1_{S_0}(c)\), \(M_f=M\), \(b_f=1\). The centered retained quadratic is
\[
 G=1_{\{T=0\}}X,\qquad X=\sum_{j\in S_0}(\xi_j^2-1).
\]
Direct calculation gives
\begin{align}
 \Var G&=2M/R,\notag\\
 \kappa_4(G)&=48M/R+12M^2(1/R-1/R^2),\label{gs:counterkappa}\\
 \E e^{tG}&=1-1/R+R^{-1}e^{-tM}(1-2t)^{-M/2},\qquad t<1/2.
 \label{gs:countermgf}
\end{align}
For fixed positive \(t\), the logarithm of \eqref{gs:countermgf} is
\(M[-t-\tfrac12\log(1-2t)]-\log R+o(1)\) as \(M\to\infty\).
Taking \(R\) large defeats any uniform bound \(CMt^2/R\) with a fixed
positive radius. This is a Gaussian linear-observation counterexample to that
inference. It uses the abstract matrix hypotheses rather than the particular
Maxwell root family; it is not a counterexample to the variance theorem or to a
separately proved Wilson estimate.
\end{proof}

There is nevertheless a valid complete mixture bound. Put
\(V_{\max}=48d_GM_f\) and \(\bar q=\E_Tv_T/V_{\max}\le C/R\), with the
zero-mark case interpreted as zero. The Gaussian determinant formula gives
\[
 \log\E_\xi e^{tG_{T,f}}
 \le\frac{\upsilon^2t^2v_T}{1-16|\upsilon t|b_f},
 \qquad16|\upsilon t|b_f<1.
\]
For \(0\le v\le V_{\max}\), convexity bounds \(e^{av}\) \emph{above} by
the chord joining its endpoint values. Averaging that bound yields
\begin{equation}
 \E_{T,\xi}e^{tG_{T,f}}
 \le1+\bar q\left[
 \exp\!\left(\frac{\upsilon^2t^2V_{\max}}
                   {1-16|\upsilon t|b_f}\right)-1\right].
 \label{gs:validmixturemgf}
\end{equation}
In particular,
\begin{equation}
 \log\E e^{tG}\le C\upsilon^2d_GM_ft^2/R,
 \qquad |\upsilon t|\sqrt{d_GM_f}\le c.
 \label{gs:smallmixtureradius}
\end{equation}
This radius depends on the whole marked mass. Gaussian second-chaos moments
also give, for \(p\ge2\),
\[
 \norm G_{L^p}
 \le C|\upsilon|p\,(\E_Tv_T^{p/2})^{1/p}
 \le C|\upsilon|p\,R^{-1/p}\sqrt{d_GM_f}.
\]
These are real mixture bounds with all offset terms retained. They provide no
volume-independent complex mixed-admission statement.

\subsection{Relation to the unresolved compact susceptibility}

The rank proof works because this specified observation is a known linear
Gaussian observation and Wick contraction gives an entrywise positive full
trace kernel. The actual compact conditional mean has neither representation
proved. Counting retained links does not estimate its nonlinear excess under
\(e^{-\beta S_{\rm out}(b)}\prod_jZ_j(b)\,\dd b\); every raw normalizer and
the full boundary and source responses remain part of that problem.

The results above therefore supply a quantitative reference benchmark, its
optimal separator exponent, a controlled shared holonomy mean, and a precise
higher-source limitation. They do not prove the required
\(\mathcal A^*\mathcal L_1^{-1}\mathcal A\le CI\), the original ordinary
heat-pressure upper curvature, or a physical-scale induction. The original
physical source factor still squares to \(c_\beta^{-2}a^{-4}\); source
normalizations are not altered by Gaussian whitening. The flat-block scale
restriction and the simultaneous-domain, current, reserve and higher-partner
obligations remain unchanged.
